\subsection{A full-AA theorem for finitely many exchangeable species}
\label{threshold_physics:species}

Let \([n]=B_1\sqcup\cdots\sqcup B_r\) be a known partition. A law is
\emph{species-exchangeable} if it is invariant under independent
coordinate permutations within every \(B_b\). Conditional on the vector
of block weights \((k_1,\ldots,k_r)\), such a law is the product
\(\bigotimes_b U_{|B_b|,k_b}\). The weights themselves can be arbitrarily
dependent.

\begin{theorem}[Simultaneous simulation on products of slices]
\label{thm:species_simulation}
Assume uniform-input acceptance simulation with a fixed polynomial query
bound \(S_0(T,1/e)\). There is a polynomial \(R(T,r,1/\epsilon)\) such
that for every known partition into \(r\) blocks and every \(T\)-query
acceptance function \(Q\), one deterministic decision tree \(D\) makes
at most \(R(T,r,1/\epsilon)\) queries and satisfies
\[
 \mathbb E_{\bigotimes_b U_{|B_b|,k_b}}|D-Q|\le\epsilon
 \quad\text{for every admissible vector }(k_1,\ldots,k_r).
\]
Consequently, for fixed \(r\), every species-exchangeable family has
full distributional AA, uniformly over its mixing law and all its model
parameters.
\end{theorem}

\begin{proof}
We extend the slice argument in Section~\ref{proof:exchangeable_queries}.
First, its biased-product lemma works for arbitrary known coordinate
biases \(q_i\in[a,1-a]\), with the same bound
\[
 S(T,e,a)=S_0(2bT,2/e),\qquad
 b=\max\left\{1,\left\lceil
 \log_2\frac{64T^2}{ae^2}\right\rceil\right\}.
\]
Indeed use a different threshold gadget of the same \(b\)-bit length at
each coordinate, rounding each \(q_i\) upward to a dyadic. Every quantum
oracle call costs at most \(2b\) uniform-bit queries. Conditional sampling
of each source block given its output bit remains independent across
coordinates, and the rounded-product coupling has conditional mismatch
at most \(2^{-b}/a\) at every coordinate. The same hybrid bound and
averaging over input-independent auxiliary seeds prove the assertion.
Replace \(S\) by a monotone polynomial majorant when convenient.

For \(T\ge1\) and \(0<\epsilon<1\), take
\[
\begin{split}
a&=(\epsilon/(64T))^2,\qquad
s=\max\{1,S(T,\epsilon/8,a)\},\\
g&=\min\{\epsilon/(16s),(\epsilon/(64T))^2\},\qquad
\eta=ag/2,\\
L&=\left\lceil4r/(\epsilon\eta^2)\right\rceil,\qquad
N_0=\left\lceil1/(a^2g^2)\right\rceil,\qquad
M=\max\{2L,L+N_0\}.
\end{split}
\]
Read every coordinate of blocks of size less than \(M\). In each
remaining block query its first \(L\) coordinates and let \(q_b\) be
their empirical density. Restrict \(Q\) to every coordinate read so far.
For a remaining block with \(q_b\le2a\), set its unread coordinates to
zero in this restricted quantum computation; with \(q_b\ge1-2a\), set
them to one. Call the resulting acceptance function on the other,
\emph{central}, coordinates \(Q_0\). Simulate \(Q_0\) under the product
law whose coordinate bias is \(q_b\) on central block \(b\), using the
vector-biased-product lemma at error \(\epsilon/8\). This is a fixed
tree for each observed transcript. If no central coordinates remain,
evaluate the fully specified computation with no further queries.
All choices depend only on the observed transcript, not on the unknown
\(k_b\). The query count is at most \(rM+s\), a polynomial in the stated
parameters.

For the analysis fix an arbitrary weight vector. After the prefixes,
the unread coordinates are independent uniform slices. Let \(p_b\)
be the unread density in block \(b\). The hypergeometric estimate from
Section~\ref{exch:sec-4} and a union bound give
\[
 \Pr[\exists b:\ |p_b-q_b|>\eta]\le r/(L\eta^2)\le\epsilon/4.
\]
On a good transcript an endpoint block has density at most \(3a\) away
from its chosen constant string. Couple the original unread input with
the input in which all endpoint blocks have been replaced by these
constant strings. Conditional on the latter input, every coordinate has
mismatch probability at most \(3a\): the endpoint slices are independent
of the central slices, and are exchangeable internally. The quantum
hybrid estimate therefore bounds the expected change of acceptance by
\(4T\sqrt{3a}<\epsilon/8\). There is no factor \(r\), since query
weights sum to one over all coordinates.

For central block \(b\), \(p_b\in[a,1-a]\) and its unread size is at
least \(N_0\). Independently couple each central slice to its
\(q_b\)-biased product as in Section~\ref{exch:sec-3}. Conditional on
the entire original central input, every coordinate has mismatch
probability at most
\[
 \eta/a+1/(2a\sqrt{N_0})\le g.
\]
The classical and quantum hybrid bounds give central error at most
\[
 \epsilon/8+sg+4T\sqrt g\le\epsilon/4.
\]
Including endpoint replacement and bad transcripts gives total error
less than \(\epsilon/4+\epsilon/8+\epsilon/4<\epsilon\), simultaneously
for every weight vector. The case \(T=0\) is immediate. Mixing the
weight vectors proves the last assertion.
\end{proof}

The fixed-species qualification matters: \(r=r_n\) produces a bound
polynomial in \(r_n\), which is not in general dimension independent.
No ability to sample the Gibbs law, determine the weight vector, or
identify a thermodynamic phase is assumed.
